\chapter{Alternating Direction Method of Multipliers}

\section{Introduction}

\section{Static Problem}

Following Bourgat, Dumay, and Glowinski~\autocite{glowinski_large_nodate}, we
model the static configuration of an inextensible flexible beam by its
centerline
\begin{equation*}
    \vec r : [0,L] \to \mathbb{R}^d, \qquad d \in \{2,3\},
\end{equation*}
where \(s\) denotes arclength along the reference centerline. The inextensibility
constraint is therefore imposed by
\begin{equation}\label{eq:addm-static-inextensibility}
    \norm{\pdv{\vec r}{s}} = 1 \quad \text{on } [0,L].
\end{equation}
This constraint is the source of the geometric nonlinearity: the bending energy
is quadratic in \(\vec r\), but the admissible set is nonconvex.

For a prescribed external force density \(\vec f(s)\), the static equilibrium
problem may be written as the constrained local minimization problem
\begin{equation}\label{eq:addm-static-minloc}
    \minloc_{\vec r \in V}
    \left\{
    \frac{EI}{2}
    \int_0^L \norm{\pdv[2]{\vec r}{s}}^2 \odif{s}
    -
    \int_0^L \vec f \cdot \vec r \odif{s}
    \right\},
\end{equation}
where \(EI>0\) is the flexural stiffness. For example, a planar gravity load can
be recovered by taking \(\vec f = -\rho g \vec e_y\), up to the sign convention
used for the vertical coordinate.

The admissible set is
\begin{equation}\label{eq:addm-static-admissible-set}
    \begin{split}
        V =
        \left\{
        \vec r \in C^1([0,L];\mathbb{R}^d)
        :
        \pdv[2]{\vec r}{s} \in L^2(0,L;\mathbb{R}^d),
        \ \norm{\pdv{\vec r}{s}} = 1,
        \ \text{and BCs hold}
        \right\}.
    \end{split}
\end{equation}
In two dimensions this reduces to the pipeline model of Bourgat, Dumay, and
Glowinski, with \(\vec r=(x,y)\) and \(x'^2+y'^2=1\). Their analysis establishes
existence of minimizers for the static problem under the endpoint conditions
considered in their work, and they also emphasize that minimizers need not be
unique. In particular, even simple endpoint constraints may admit several
equilibrium branches. Since the present work is primarily concerned with the
numerical enforcement of inextensibility, we use those analytical results as
motivation for the constrained variational formulation rather than reproducing
the full bifurcation analysis.

Introducing a Lagrange multiplier \(\lambda(s)\) for
\eqref{eq:addm-static-inextensibility} gives the formal Euler--Lagrange
equation
\begin{equation}\label{eq:addm-static-strong-form}
    EI \pdv[4]{\vec r}{s}
    -
    \pdv*{\left(\lambda \pdv{\vec r}{s}\right)}{s}
    =
    \vec f,
    \qquad
    \norm{\pdv{\vec r}{s}} = 1,
\end{equation}
with boundary conditions determined by the chosen support constraints. This
strong form is useful for interpretation, but it is not the most convenient
starting point for computation because the multiplier is coupled directly to the
nonlinear arclength constraint.

The key computational idea in Bourgat, Dumay, and Glowinski is to split the
derivative and constraint by introducing an auxiliary tangent variable
\begin{equation}\label{eq:addm-static-split-variable}
    \vec p = \pdv{\vec r}{s},
    \qquad
    \norm{\vec p}=1.
\end{equation}
The nonlinear constraint is then isolated as a pointwise projection onto the unit
sphere in \(\mathbb{R}^d\), while the relation
\(\vec p-\pdv{\vec r}{s}=0\) is a linear constraint. An augmented Lagrangian can
therefore be applied to this linear constraint:
\begin{equation}\label{eq:addm-static-augmented-lagrangian}
    \begin{split}
        \mathcal{L}_r(\vec r,\vec p,\vec \eta)
         & =
        \frac{EI}{2}
        \int_0^L \norm{\pdv[2]{\vec r}{s}}^2 \odif{s}
        -
        \int_0^L \vec f \cdot \vec r \odif{s}
        \\                                & \qquad +
                                         \int_0^L \vec \eta \cdot
                                         \left(\vec p-\pdv{\vec r}{s}\right)\odif{s}
        +
                                         \frac{r}{2}
        \int_0^L
                                         \norm{\vec p-\pdv{\vec r}{s}}^2 \odif{s},
    \end{split}
\end{equation}
where \(r>0\) is the penalty parameter and
\(\vec \eta\) is the multiplier associated with the split constraint. Alternating
minimization with respect to \(\vec p\) and \(\vec r\), followed by a multiplier
update, yields the alternating-direction method used to enforce
inextensibility. The \(\vec p\)-subproblem is local in \(s\), while the
\(\vec r\)-subproblem is a linear fourth-order beam solve with the current
multiplier and tangent field fixed.

\subsection{Discrete Static Problem}

Let
\begin{equation*}
    0 = s_0 < \cdots < s_N = L,
    \qquad
    h = \max_{0 \leq i < N}(s_{i+1}-s_i),
\end{equation*}
be a partition of the centerline. Since the static energy contains second
derivatives of \(\vec r\), the natural conforming approximation space is a
\(C^1\) cubic Hermite space,
\begin{equation}\label{eq:addm-hermite-space}
    V_h =
    \left\{
    \vec r_h \in C^1([0,L];\mathbb{R}^d)
    :
    \vec r_h|_{[s_i,s_{i+1}]}
    \in \mathbb{P}_3([s_i,s_{i+1}];\mathbb{R}^d),
    \ i=0,\ldots,N-1
    \right\}.
\end{equation}
The degrees of freedom are the nodal positions and tangents,
\begin{equation*}
    \left\{
    \vec r_h(s_i),
    \pdv{\vec r_h}{s}(s_i)
    \right\}_{i=0}^N.
\end{equation*}
This is the same approximation used by Bourgat, Dumay, and Glowinski in the
planar case. It is well suited to the beam energy because
\(V_h \subset H^2(0,L;\mathbb{R}^d)\), so the curvature term is defined without
introducing additional weak variables.

The finite dimensional admissible set is obtained by imposing the boundary
conditions together with finitely many tangent-length constraints. The basic
choice is the nodal constraint set
\begin{equation}\label{eq:addm-discrete-admissible-set}
    V_h^{\mathrm{inext}}
    =
    \left\{
    \vec r_h \in V_h
    :
    \norm{\pdv{\vec r_h}{s}(s_i)} = 1,
    \ i=0,\ldots,N,
    \ \text{and the BCs hold}
    \right\}.
\end{equation}
The corresponding discrete static problem is
\begin{equation}\label{eq:addm-discrete-static-problem}
    \minloc_{\vec r_h \in V_h^{\mathrm{inext}}}
    J_h(\vec r_h),
\end{equation}
where
\begin{equation}\label{eq:addm-discrete-static-energy}
    J_h(\vec r_h)
    =
    \frac{EI}{2}
    \int_0^L \norm{\pdv[2]{\vec r_h}{s}}^2 \odif{s}
    -
    \int_0^L \vec f \cdot \vec r_h \odif{s}.
\end{equation}
Thus the static finite element problem is a finite dimensional nonlinear
programming problem with quadratic equality constraints. The objective is
quadratic in the unknown centerline degrees of freedom, while the feasible set
is nonconvex because of the unit-tangent constraints.

Bourgat, Dumay, and Glowinski also describe a stronger pointwise approximation
of the inextensibility constraint. In addition to the nodal conditions in
\eqref{eq:addm-discrete-admissible-set}, one may enforce
\begin{equation}\label{eq:addm-discrete-midpoint-constraint}
    \norm{\pdv{\vec r_h}{s}(s_{i+1/2})} = 1,
    \qquad
    s_{i+1/2} = \frac{s_i+s_{i+1}}{2},
    \qquad
    i=0,\ldots,N-1.
\end{equation}
This nodal-plus-midpoint constraint controls the tangent length inside each
element, rather than only at the endpoints. It is therefore more accurate in
configurations with strong curvature variation, although the convergence proof
in~\autocite{glowinski_large_nodate} is given for the nodal constraint version.

The convergence result used here can be summarized as follows. Suppose the
continuous admissible set is nonempty, the boundary data are compatible with an
arclength-parametrized curve, and the discrete admissible sets
\(V_h^{\mathrm{inext}}\) are nonempty for sufficiently small \(h\). For the
planar problem with cubic Hermite approximation and nodal tangent constraints,
Bourgat, Dumay, and Glowinski prove that isolated local minimizers of
\eqref{eq:addm-static-minloc} are approximated by local minimizers of
\eqref{eq:addm-discrete-static-problem}. More precisely, if \(\vec r\) is an
isolated local minimizer, then for sufficiently small \(h\) there exists a
discrete local minimizer \(\vec r_h\) such that
\begin{equation}\label{eq:addm-discrete-local-convergence}
    \vec r_h \to \vec r
    \qquad
    \text{strongly in } H^2(0,L;\mathbb{R}^d)
    \quad \text{as } h \to 0.
\end{equation}
They also prove a global version: any family of discrete global minimizers has a
subsequence converging strongly in \(H^2\) to a global minimizer of the
continuous problem.

The proof relies on three standard ingredients. First, the Hermite interpolant
is consistent in \(H^2\). Second, bounded discrete energies give compactness of
the centerline sequence in \(H^2\), with strong convergence in \(C^1\) along
subsequences. Third, the nodal tangent constraints pass to the limit, so the
limit curve satisfies \(\norm{\vec r'}=1\). Weak lower semicontinuity of the
bending energy then identifies the limit as a minimizer. The result is
therefore a convergence statement for approximating a selected equilibrium
branch; it does not imply uniqueness of the continuous or discrete static
equilibrium.

The analysis in~\autocite{glowinski_large_nodate} is written for \(\vec
r=(x,y)\in\mathbb{R}^2\). The formulation above is kept in \(\mathbb{R}^d\),
\(d\in\{2,3\}\), because the same compactness and consistency mechanism applies
componentwise to the vector-valued Hermite approximation. The three-dimensional
case should therefore be understood as the corresponding vector generalization
of the discrete model, while the cited theorem is the planar result.

\subsection{Alternating-Direction Static Solve}

The split augmented Lagrangian in \eqref{eq:addm-static-augmented-lagrangian}
leads to a simple alternating iteration. Let \(V_h^{\mathrm{bc}}\) denote the
affine Hermite space satisfying the prescribed boundary conditions, and let
\(W_h^0\) be the corresponding space of homogeneous admissible variations.
Given \(\vec r_h^n\) and \(\vec \eta_h^n\), the first step updates the tangent
field by minimizing with respect to \(\vec p\) while holding the centerline
fixed:
\begin{equation}\label{eq:addm-p-update-min}
    \vec p_h^{n+1}(s)
    =
    \operatorname*{arg\,min}_{\norm{\vec p}=1}
    \left\{
    \vec \eta_h^n(s)\cdot \vec p
    +
    \frac{r}{2}
    \norm{
        \vec p
        -
        \pdv{\vec r_h^n}{s}(s)
    }^2
    \right\}.
\end{equation}
This is the vector form of the pointwise minimization problem appearing in
Bourgat, Dumay, and Glowinski. Since \(\norm{\vec p}=1\), the quadratic term in
\(\vec p\) is constant on the constraint set except for its linear coupling with
\(\pdv{\vec r_h^n}{s}\). Defining
\begin{equation}\label{eq:addm-projection-vector}
    \vec z_h^n(s)
    =
    \vec \eta_h^n(s)
    -
    r\pdv{\vec r_h^n}{s}(s),
\end{equation}
the update is the projection formula
\begin{equation}\label{eq:addm-p-update-projection}
    \vec p_h^{n+1}(s)
    =
    -
    \frac{\vec z_h^n(s)}{\norm{\vec z_h^n(s)}}
    =
    \frac{
        r\pdv{\vec r_h^n}{s}(s)-\vec \eta_h^n(s)
    }{
        \norm{
            r\pdv{\vec r_h^n}{s}(s)-\vec \eta_h^n(s)
        }
    },
    \qquad
    \vec z_h^n(s)\neq \vec 0.
\end{equation}
If \(\vec z_h^n(s)=\vec 0\), every unit vector minimizes
\eqref{eq:addm-p-update-min} at that point. The original article notes that this
degenerate case is avoided in practice by taking the penalty parameter \(r\)
sufficiently large.

With \(\vec p_h^{n+1}\) fixed, the centerline update is the unconstrained
quadratic minimization
\begin{equation}\label{eq:addm-r-update-min}
    \vec r_h^{n+1}
    =
    \operatorname*{arg\,min}_{\vec v_h\in V_h^{\mathrm{bc}}}
    \mathcal{L}_r(\vec v_h,\vec p_h^{n+1},\vec \eta_h^n).
\end{equation}
Equivalently, \(\vec r_h^{n+1}\in V_h^{\mathrm{bc}}\) satisfies
\begin{equation}\label{eq:addm-r-update-weak}
    \begin{split}
        EI
        \int_0^L
        \pdv[2]{\vec r_h^{n+1}}{s}
        \cdot
        \pdv[2]{\vec w_h}{s}
        \odif{s}
        +
        r
        \int_0^L
        \pdv{\vec r_h^{n+1}}{s}
        \cdot
        \pdv{\vec w_h}{s}
        \odif{s}
         & =
        \int_0^L
        \vec f \cdot \vec w_h
        \odif{s}
        \\
         & \quad +
        \int_0^L
        \left(
        \vec \eta_h^n
        +
        r\vec p_h^{n+1}
        \right)
        \cdot
        \pdv{\vec w_h}{s}
        \odif{s},
    \end{split}
\end{equation}
for all \(\vec w_h\in W_h^0\). The multiplier is then advanced by
\begin{equation}\label{eq:addm-multiplier-update}
    \vec \eta_h^{n+1}
    =
    \vec \eta_h^n
    +
    \alpha
    \left(
    \vec p_h^{n+1}
    -
    \pdv{\vec r_h^{n+1}}{s}
    \right),
\end{equation}
where \(\alpha>0\) is a relaxation parameter. The choice \(\alpha=r\) is the
standard one motivated by the Douglas--Rachford interpretation of the
alternating-direction method.

The main computational advantage is that the two subproblems have very
different structures. The \(\vec p_h\)-update is local in \(s\): after
discretization it is only a normalization of the vector \(r\pdv{\vec
r_h^n}{s}-\vec \eta_h^n\) at the chosen constraint or quadrature points. The
\(\vec r_h\)-update is global, but it is linear. In matrix form its left-hand
side is
\begin{equation}\label{eq:addm-linear-system-matrix}
    \mat A
    =
    EI\mat K_b
    +
    r\mat K_s,
\end{equation}
where \(\mat K_b\) is the Hermite bending stiffness matrix associated with
\(\int_0^L \vec r_h''\cdot\vec w_h''\odif{s}\), and \(\mat K_s\) is the
first-derivative matrix associated with
\(\int_0^L \vec r_h'\cdot\vec w_h'\odif{s}\). For fixed \(r\), this matrix does
not change with the nonlinear iteration. With essential boundary conditions
that remove rigid translations and other null modes, \(\mat A\) is sparse,
symmetric, and positive definite. The same scalar Hermite matrix is used
componentwise in \(\mathbb{R}^d\), so the vector problem has a block diagonal
structure and a single factorization can be reused for all components and all
alternating-direction iterations.

\section{Dynamic Problem}

\subsection{\texorpdfstring{Newmark-\(\beta\)}{Newmark-beta}}

For a general second-order in time dynamical system
\begin{equation*}
    \mat M \ddot{\vec u} + \mat C \dot{\vec u} + f^{\rm int}(\vec u) = f^{\rm ext}
\end{equation*}
where \(\mat M\) is a mass matrix, and \(\mat C\) is a damping matrix. Newmark-\(\beta\) scheme provides for the update
\begin{align*}
    \vec u_{n+1}        & = \vec u_n + \Delta t \dot{\vec u}_n + \frac12 \Delta t^2 \ddot {\vec u}_\beta \\
    \dot {\vec u}_{n+1} & = \dot{\vec u}_n + \Delta t \dot{\vec u}_\gamma
\end{align*}
where the acceleration vector \(\ddot{\vec u}_\beta\) and the velocity vector \(\ddot{\vec u}_\gamma\) are defined as
\begin{align*}
    \ddot{\vec u}_\beta   & = (1-2\beta) \ddot{\vec u}_n + 2\beta\ddot{\vec u}_{n+1}  \\
    \dot{\vec u}_{\gamma} & = (1- \gamma)\ddot{\vec u}_n + \gamma \ddot{\vec u}_{n+1}
\end{align*}
Substituted into the updates we have
\begin{align*}
    \vec u_{n+1}        & = \vec u_n + \Delta t \dot{\vec u}_n + \frac12 \Delta t^2  \left[(1-2\beta) \ddot{\vec u}_n + 2\beta \ddot{\vec u}_{n+1}\right] \\
    \dot {\vec u}_{n+1} & = \dot{\vec u}_n + \Delta t \left[(1- \gamma)\ddot{\vec u}_n + \gamma \ddot{\vec u}_{n+1}\right]
\end{align*}
which gives us a stencils for \(\ddot{\vec u}_{n+1}\) and \(\dot{\vec u}_{n+1}\) which are
\begin{align*}
    \ddot{\vec u}_{n+1} & = \frac{1}{\beta \Delta t^2}(\vec u_{n+1} - \vec{u}_n - \Delta t \dot{\vec u}_n) - \frac{1-2\beta}{2\beta}\ddot{\vec u}_n \\
    \dot{\vec u}_{n+1}  & = \dot{\vec u}_n + \Delta t[(1-\gamma)\ddot{\vec u}_n + \gamma \ddot{\vec u}_{n+1}]
\end{align*}
In our time dependent E-B problem we have
\begin{equation*}
    \mu \ddot{\vec r}(s) + \underbrace{EI \vec{r}^{(4)}(s) - (\lambda(s) \vec r'(s))'}_{\vec f^{\rm int}} =  \underbrace{\vec F(s,t)}_{\vec f^{\rm ext}}
\end{equation*}
which has no damping. Therefore, we only need to discretize \(\ddot{\vec r}^{n+1}\):
\begin{equation*}
\end{equation*}

\section{Validation}

\subsection{Elliptic Integral Solution}

We use a static equilibrium solution for the tip of a cantilevered inextensible
beam given by Bisshopp and Drucker \cite{bisshopp_large_1945}. The solution
gives the finite-deflection tip position of an Euler--Bernoulli beam of length
\(L\), flexural rigidity \(B=EI\), and vertical tip load \(P\). It preserves
arc length and evaluates the curvature with respect to the deformed centerline,
so the result is appropriate for large rotations of an inextensible beam. In
the validation problem we use the solution only as a reference for the vertical
tip drop \(\delta\) and horizontal tip retreat \(A=L-x_{\rm tip}\).

Define the nondimensional load
\begin{equation*}
    a = \sqrt{\frac{P L^2}{B}} .
\end{equation*}
The elliptic modulus \(k\in[1/\sqrt{2},1)\) is determined from the scalar
equation
\begin{equation*}
    a = K(k) - F(\theta_1 \mid k),
    \qquad
    \theta_1 =
    \sin^{-1}\left(\frac{1}{\sqrt{2}k}\right),
\end{equation*}
where \(K(k)\) is the complete elliptic integral of the first kind and
\(F(\theta\mid k)\) is the incomplete elliptic integral of the first kind. In
the benchmark this equation is solved by bisection for \(k\).

Once \(k\) is known, the tip displacement follows from elliptic integrals of
the second kind:
\begin{align*}
    \frac{\delta}{L}
     & = 1 -
    \frac{2\left(E(k)-E(\theta_1\mid k)\right)}{a}, \\
    \frac{x_{\rm tip}}{L}
     & = \frac{\sqrt{2}\sqrt{2k^2-1}}{a}.
\end{align*}
Thus the horizontal retreat used for comparison is
\begin{equation*}
    A
    = L - x_{\rm tip}
    = L\left(1-\frac{\sqrt{2}\sqrt{2k^2-1}}{a}\right).
\end{equation*}
For \(P=0\), the limiting values are \(k=1/\sqrt{2}\), \(\delta=0\), and
\(A=0\). The numerical tests compare the computed tip coordinates with
\(\abs{y_{\rm tip}}\approx\delta\) and
\(\abs{L-x_{\rm tip}}\approx A\).

\begin{figure}[H]
    \centering
    \import{figure/vector}{bisshopp-drucker.tex}
    \caption{Non-dimensional tip deflection \(y_{\rm tip}/L\) and tip retreat \(x_{\rm tip}/L\) with respect to non-dimensional tip load \(PL^2/B\).}
    \label{fig:bisshopp-drucker-reference}
\end{figure}

The static ADDM discretization is tested against this reference solution by
solving the same cantilever problem on a sequence of uniform meshes. The beam
has \(L=1\), \(EI=1\), and a downward tip load \(P=1\). For each mesh, the ADDM
solve produces a numerical tip position \(\vec r_{h,\rm tip}\). This is
compared with the Bisshopp--Drucker tip position \(\vec r_{\rm
tip}=(L-A,-\delta)\), where \(A\) and \(\delta\) are computed from the
elliptic-integral solution above. The scalar error reported in
Figure~\ref{fig:bisshopp-drucker-convergence} is
\begin{equation*}
    \norm{\vec \epsilon_{\rm tip}}
    =
    \left[
        \left(\abs{L-x_{h,\rm tip}}-A\right)^2
        +
        \left(\abs{y_{h,\rm tip}}-\delta\right)^2
        \right]^{1/2}.
\end{equation*}
The test is carried out for the node counts
\(11,22,44,88,\) and \(172\), and the error is plotted against
\(1/\Delta s\). The error decreases monotonically under mesh refinement. The
log-log slope is steeper than two over this range, which confirms the expected
second-order-in-space convergence of the static ADDM discretization for this
tip quantity.

\begin{figure}[H]
    \centering
    \import{figure/vector}{addm-static-convergence.tex}
    \caption{Log-log plot of tip error \(\norm{\vec {\epsilon}_{\rm tip}}\) versus \(1/\Delta s\) of the beam mesh.}
    \label{fig:bisshopp-drucker-convergence}
\end{figure}

\subsection{Pinned Manufactured Solution}

To test the temporal convergence of the Newmark integrator using the method of
manufactured solutions. The manufactured motion is chosen to satisfy the
inextensibility constraint exactly, so that the observed error is dominated by
the time discretization and by the nonlinear solve tolerance rather than by a
constraint-inconsistent reference trajectory.

The reference solution used in the benchmark is the shifted circular arc
implemented in \texttt{ManufacturedDynamicResult1}. We take
\begin{equation*}
    \vec r_{\rm ex}(s,t)
    =
    \frac1{\phi(t)}
    \begin{bmatrix}
        \cos(s\phi(t)) - 1 \\
        \sin(s\phi(t))
    \end{bmatrix},
    \qquad
    \phi(t)=e^t,
    \qquad
    s\in[0,\pi/2].
\end{equation*}
The implementation uses this as a planar solution with zero third component.
The shift in the first component pins the left endpoint,
\(\vec r_{\rm ex}(0,t)=\vec 0\), but does not affect any spatial derivative.

Let \(\theta(s,t) = s\phi(t)\),
\begin{equation*}
    \vec e_r = \begin{bmatrix}\cos{\theta} \\ \sin{\theta} \end{bmatrix}, \qquad
    \vec e_\theta = \begin{bmatrix}-\sin{\theta} \\ \cos{\theta}\end{bmatrix}, \qquad
    \vec e_x = \begin{bmatrix}1 \\ 0\end{bmatrix}.
\end{equation*}
Then the manufactured solution can be written as
\begin{equation*}
    \vec r_{\rm ex}(s,t) = \frac1\phi \left(\vec e_r - \vec e_x\right).
\end{equation*}
The relevant spatial derivatives of the moving basis are
\begin{equation*}
    \pdv{\vec e_r}{s} = \phi \vec e_\theta, \qquad
    \pdv{\vec e_\theta}{s} = -\phi \vec e_r.
\end{equation*}
Therefore
\begin{equation*}
    \pdv{\vec r_{\rm ex}}{s} = \vec e_\theta,
    \qquad
    \pdv[2]{\vec r_{\rm ex}}{s} = -\phi\vec e_r,
    \qquad
    \pdv[3]{\vec r_{\rm ex}}{s} = -\phi^2\vec e_\theta,
    \qquad
    \pdv[4]{\vec r_{\rm ex}}{s} = \phi^3\vec e_r.
\end{equation*}
In particular,
\begin{equation*}
    \norm{\pdv{\vec r_{\rm ex}}{s}} = 1.
\end{equation*}
Thus the reference motion is exactly inextensible for any sufficiently smooth
choice of \(\phi\). The bending term is
\begin{equation*}
    EI \vec r_{\rm ex}^{(4)} = EI \phi^3 \vec e_r.
\end{equation*}

For the inertial term, the time-dependence of the endpoint shift must be
retained. Since
\begin{equation*}
    \pdv{\vec e_r}{t} = s\dot{\phi}\vec e_\theta, \qquad
    \pdv{\vec e_\theta}{t} = -s\dot{\phi}\vec e_r.
\end{equation*}
Differentiating once in time gives
\begin{equation*}
    \pdv{\vec r_{\rm ex}}{t}
    =
    -\frac{\dot{\phi}}{\phi^2}\left(\vec e_r-\vec e_x\right)
    + \frac{s\dot{\phi}}{\phi}\vec e_\theta.
\end{equation*}
Differentiating again gives
\begin{equation*}
    \pdv[2]{\vec r_{\rm ex}}{t}
    =
    \left(-\frac{\ddot{\phi}}{\phi^2}
    + \frac{2\dot{\phi}^2}{\phi^3}
    - \frac{s^2\dot{\phi}^2}{\phi}\right)\vec e_r
    +
    \left(\frac{\ddot{\phi}}{\phi^2}
    - \frac{2\dot{\phi}^2}{\phi^3}\right)\vec e_x
    +
    \left(\frac{s\ddot{\phi}}{\phi}
    - \frac{2s\dot{\phi}^2}{\phi^2}\right)\vec e_\theta.
\end{equation*}
For the choice \(\phi(t)=\exp{(t)}\), so that \(\dot{\phi}=\ddot{\phi}=\phi\), this reduces to
\begin{equation*}
    \ddot{\vec r}_{\rm ex}
    =
    \left(\frac1\phi - s^2\phi\right)\vec e_r
    - \frac1\phi \vec e_x
    - s\vec e_\theta.
\end{equation*}
This is the acceleration used by the reference solution, written in the moving
basis.

For the linear dynamic Euler--Bernoulli equation
\begin{equation*}
    \rho \ddot{\vec r} + EI \vec r^{(4)} = \vec f
\end{equation*}
the corresponding manufactured load is
\begin{equation*}
    \vec f_{\rm linear}
    =
    \rho\left[
        \left(\frac1\phi - s^2\phi\right)\vec e_r
        - \frac1\phi \vec e_x
        - s\vec e_\theta
        \right]
    + EI\phi^3\vec e_r.
\end{equation*}
Equivalently, this can be written as
\begin{equation*}
    \vec f_{\rm linear}
    =
    \left[\rho\left(\frac1\phi - s^2\phi\right)+EI\phi^3\right]\vec e_r
    - \frac{\rho}{\phi}\vec e_x
    - \rho s\vec e_\theta.
\end{equation*}

For the nonlinear inextensible benchmark, the load convention used in the
reference implementation is
\begin{equation*}
    \vec f_{\rm nonlinear}
    =
    \rho\ddot{\vec r}_{\rm ex}
    + EI\vec r_{\rm ex}^{(4)}
    + \pdv*{\left(\lambda \pdv{\vec r_{\rm ex}}{s}\right)}{s}.
\end{equation*}
The multiplier contribution is
\begin{equation*}
    \pdv*{\left(\lambda \pdv{\vec r_{\rm ex}}{s}\right)}{s}
    =
    \pdv{\lambda}{s}\vec e_\theta - \lambda\phi\vec e_r.
\end{equation*}
With the manufactured choice \(\lambda=1\), this contribution is
\(-\phi\vec e_r\), and hence
\begin{equation*}
    \vec f_{\rm nonlinear}
    =
    \vec f_{\rm linear} - \phi\vec e_r,
\end{equation*}
or equivalently
\begin{equation*}
    \vec f_{\rm nonlinear}
    =
    \left[\rho\left(\frac1\phi - s^2\phi\right)+EI\phi^3-\phi\right]\vec e_r
    - \frac{\rho}{\phi}\vec e_x
    - \rho s\vec e_\theta.
\end{equation*}

\subsection{Linear Convergence Study}
In the linear study, the multiplier field is not solved for explicitly. The
boundary values of \(\vec r\) and \(\pdv{\vec r}{s}\) are prescribed from
\(\vec r_{\rm ex}\), and the body force is set to \(\vec f_{\rm linear}\). We
fix the spatial resolution at \(\Delta s=(\pi/2)/350\) and vary the time step
over
\begin{equation*}
    \log{1/\Delta t} \in \{1.3,1.4,1.5,1.6,1.7,1.8,1.9,2,2.1,2.2\}.
\end{equation*}
At the final time \(T=1\), the error is measured against the manufactured
solution using the discrete beam norm
\begin{equation*}
    E_h(\Delta t) = \norm{\vec r_h(\cdot,T) - \vec r_{\rm ex}(\cdot,T)}_2.
\end{equation*}
The spatial grid is kept fixed and sufficiently fine so that the measured slope
primarily reflects temporal error.

\begin{figure}[H]
    \centering
    \import{figure/vector}{addm-linear-convergence.tex}
    \caption{The convergence rate for linear case.}
    \label{fig:newmark-addm-linear-convergence}
\end{figure}

Figure~\ref{fig:newmark-addm-linear-convergence} shows the expected monotone
decrease in error as the time step is refined. On the logarithmic axes, the
points should lie close to a straight line if the asymptotic temporal regime
has been reached. For the average-acceleration Newmark method and a smooth
manufactured load, the expected temporal order is second order. The final
discussion should compare the fitted slope against this expected value and
identify whether the coarsest or finest time steps should be excluded from the
fit.

% TODO: Insert the fitted linear-study slope and the exact data range used for
% the regression. TODO: Check whether the finest time step is affected by the
% fixed spatial resolution or roundoff before reporting the final observed
% order.

\subsection{Nonlinear Convergence Study}

In the nonlinear study, the same manufactured trajectory is used, but the
multiplier field is solved simultaneously with the beam position. The external
forcing is set to \(\vec f_{\rm nonlinear}\), corresponding to the manufactured
choice \(\lambda=1\). This test exercises both the Newmark discretization and
the nonlinear ADDM iteration used to enforce inextensibility.

The time-step sequence and final-time error measurement are the same as in the
linear study. The nonlinear tolerances used are \(\epsilon =
\{10^{-7},10^{-8},10^{-9}\}\), along with \(r=10^3\).

% TODO: Record the nonlinear tolerance, relaxation parameter, and maximum outer
% iteration count used for the finalized run. TODO: Decide whether to report a
% constraint error alongside the displacement error. 

\begin{figure}[H]
    \centering
    \import{figure/vector}{addm-nonlinear-convergence.tex}
    \caption{Convergence of the nonlinear case.}
    \label{fig:newmark-addm-nonlinear-convergence}
\end{figure}

% TODO: Insert the fitted nonlinear-study slope and compare it directly with the
% linear-study slope. TODO: Discuss any plateau, non-monotone point, or outlier
% once the finalized convergence data are fixed.
